\paragraph{Part (ii).}
As $\lambda\to\infty$, we have $y\to0$, $\theta\to1$, $p\to1/k$, $w\to1/(k-1)$, and $u\to k^2/(k-1)$, and
$D$
tends to the unit point mass at $q=k$, since $r_i\to1$ as $x\to\infty$. In~\eqref{eq:torres-identidad},
after
cancellation of $\frac12\log y$, we are left with
\[
\begin{aligned}
 R'_k(\infty)={}&\log k!-k\log\frac{k^2}{k-1}-k^2\log\bigl(1-\tfrac1k\bigr)
 +\tfrac12(k-1)^2\log\bigl(1-(k-1)^{-2}\bigr)\\
 &+\tfrac12\log\Bigl(1+\frac{k-2}k\Bigr)-\log\bigl(1-\tfrac1k\bigr)
 +(k-1)\log\Bigl(1+\frac1{k-1}\Bigr)-\tfrac12\log(4\pi k).
\end{aligned}
\]
Since $-k\log(k^2/(k-1))-k^2\log(1-1/k)=-k\log k-k(k-1)\log(1-1/k)$ and
$\frac12\log((2k-2)/k)-\log((k-1)/k)=\frac12\log(2k/(k-1))$, this gives the stated expression. For the
expansion we use Stirling's formula, which has no term of order $k^{-2}$, together with the expansions
below.
\[
\begin{aligned}
 \log k!&=k\log k-k+\tfrac12\log(2\pi k)+\tfrac1{12k}+O(k^{-3}),\\
 -k(k-1)\log\bigl(1-\tfrac1k\bigr)&=(k-1)\sum_{m\ge1}\frac1{m\,k^{m-1}}
 =k-\tfrac12-\tfrac1{6k}-\tfrac1{12k^2}+O(k^{-3}),\\
 \tfrac12(k-1)^2\log\bigl(1-(k-1)^{-2}\bigr)&=-\tfrac12-\tfrac1{4(k-1)^2}+O(k^{-4})
 =-\tfrac12-\tfrac1{4k^2}+O(k^{-3}),\\
 \tfrac12\log\frac{2k}{k-1}&=\tfrac12\log2+\tfrac1{2k}+\tfrac1{4k^2}+O(k^{-3}),\\
 (k-1)\log\frac k{k-1}&=1-\tfrac1{2k}-\tfrac1{6k^2}+O(k^{-3}).
\end{aligned}
\]
On summing, we have $\frac12\log(2\pi k)+\frac12\log2-\frac12\log(4\pi k)=0$ and $-\frac12-\frac12+1=0$. The
coefficient
of $1/k$ is $\frac1{12}-\frac16+\frac12-\frac12=-\frac1{12}$, and that of $1/k^2$ is
$-\frac1{12}-\frac14+\frac14-\frac16=-\frac14$. Hence $R'_k(\infty)=-\frac1{12k}-\frac1{4k^2}+O(k^{-3})$.

\paragraph{Part (iii).}
With $k$ fixed and $\lambda\to0^+$, we have $y\to1$ and $\theta\to0$. By~\eqref{eq:torres-punto},
$p/\lambda\to1$, so $\theta=kp\sim k\lambda$. Then $\log Z_k\to0$. Moreover,
$\log Z_B=-2k\log y+k^2\log(1-p)\to0$, $\det(I+M)\to1$, and $\mathsf P\to1$, since the $q=0$ term in
$\sum_q\beta(q)^2$ is $y^{2k}\to1$. On the other hand,
$\frac12\log(4\pi ky\theta)=\frac12\log(4\pi k^2\lambda)+o(1)\to-\infty$. Hence
$R'_k=-\frac12\log(4\pi k^2\lambda)+o(1)$.

\begin{observacion}[Interpretation]
When there are very few rooks, the Gaussian $P_G$ for the balance $N_R-N_C$ is very narrow because its
variance
$2ky\theta$ tends to zero, but the determinant has no small factor there to compensate for it. The
zero-mode factor overcorrects. In the dense regime the factor of $\det(I+M)$ that tends to zero
is proportional to $y$ (Lemma~\ref{lem:torres-lazo}) and cancels against the variance.

The argument uses the fact that the number of perfect matchings between $q$ occupied rows and $q$ occupied
columns is $q!$, which makes the average in Lemma~\ref{lem:torres-identidad} explicit. For queens, the
corresponding factor is the number of placements that also respect the diagonal constraints. Counting them
is the
entire problem, so the argument of this appendix does not carry over to queens.
\end{observacion}
